\subsection{Data Collection and Political-Corruption Classifier}
\label{sec:method-political-classifier}

News articles were collected as part of the RESPOND multilingual news corpus. The data-collection workflow queried the News API using country-specific corruption keyword dictionaries with declensions and morphological variants. The collection code and keyword-generation notebooks are maintained in the related RESPOND media repository.\footnote{\url{https://github.com/annekroon/RESPOND_media/tree/main/data-collection/news-collection/news-api}} The present repository starts from the exported country-level article files and implements the cleaning, deduplication, classification, and validation workflow used to identify articles primarily about political corruption.

The first stage cleaned and deduplicated the country-level news files. Article text was standardized by using the combined title/body field available in the exported data and storing it as \texttt{article\_text}. Text normalization removed excess whitespace and repaired common character-encoding artefacts. Rows marked as duplicates by the source data were removed, articles with missing or very short text were excluded, and exact and near-exact duplicates were removed using normalized text hashes. The resulting cleaned corpus contained 2,264,916 articles from Bulgaria, France, Hungary, Italy, the Netherlands, Serbia, Sweden, Ukraine, and the United Kingdom.

The second stage trained a classifier to distinguish articles primarily about political corruption from articles that mention corruption only incidentally or concern non-political misconduct. Political corruption was defined as the misuse of public or political office, authority, or resources for private or political gain. The positive class includes articles focused on allegations, investigations, charges, or proven instances of corruption by political officials. Articles focused on private fraud, ordinary crime, or misconduct by non-political actors were treated as negative cases.

Because manual labels were limited, we used a silver-labeling strategy. Targeted article batches were selected from the cleaned corpus, prioritizing uncertain cases, predicted positives, predicted negatives, and countries with weaker validation performance. These batches were sent to the UvA LLM proxy for English translation and structured label suggestions. The LLM labels were treated as silver labels for model training rather than as final human annotations. A manually annotated validation set was kept separate for evaluation. After an initial country-level validation check showed low positive support for the United Kingdom, a UK-focused calibration batch was created and labelled. A smaller set of UK cases was then manually reviewed and added to the validation benchmark as a UK validation supplement.

Candidate classifiers were evaluated against the human validation benchmark. The final selected classifier used \texttt{intfloat/multilingual-e5-large} sentence embeddings and a class-balanced logistic regression model trained on the combined silver-label batches 1 and 2. The UK calibration batch was evaluated as a sensitivity check but was not used in the main classifier because the gain over the original combined silver-label model was modest. The decision threshold was set to 0.40 based on validation-set threshold sweeps. The final validation set contained 502 manually labelled articles, including 141 political-corruption articles. On this validation set, the final classifier achieved accuracy = 0.863, political-corruption precision = 0.731, political-corruption recall = 0.809, political-corruption F1 = 0.768, macro F1 = 0.835, and weighted F1 = 0.865.

The final classifier was applied to all cleaned country files. Articles with predicted probabilities at or above the 0.40 threshold were retained as the political-corruption corpus for subsequent analyses of victim visibility, corruption domain, case scope, and attention over time. For descriptive attention analyses, the numerator is the number of articles classified as political corruption in a country-period, and the denominator is total news coverage in that country-period rather than the initial corruption-query corpus.
